\section{Finding the cheapest successful audit}
\begin{frame}[t,label=solver]{The algorithm}
\small
\begin{enumerate}
\setlength{\itemsep}{9pt}
\item Find the cheapest plan that supports a full-market stable matching,
subject to constraints learned so far. This is the \emph{master problem}.
\item Run DA on that plan. Check for blocking pairs in the full market.
\item If a pair blocks, add a constraint excluding that failure. Solve again.
\item Stop when DA passes, or the lower bound reaches the cost of a
previously verified plan.
\end{enumerate}
\sources{Logic-based Benders decomposition (Hooker and Ottosson, 2003), using the
stable-matching polytope (Rothblum, 1992). Equations: Appendix 2.}
\end{frame}

\begin{frame}[t,label=prefix]{The part of each list that matters}
\small
Suppose a buyer ranks $x\succ y\succ z$ and DA assigns it to $y$.
Its \emph{prefix} contains $x$ and $y$, even if $x$ was hidden.

For a visible graph $H\subseteq E$, let $M=\DA(H)$. The full prefixes are
\[
 \Pi(M)=\{bs\in E:M(b)=\bot\ \text{or}\ s\succeq_b M(b)\}.
\]
An unmatched buyer contributes its whole list.

\medskip
DA never proposes below its eventual partner. Keep every prefix pair
visible or hidden exactly as before, and the same DA run replays.

\medskip
To fix a failed outcome, at least one of these visibility choices must change.
\sources{Proof: Appendix 3. Related DA-prefix reasoning appears in Aguiar and Hong (2026, preprint).}
\end{frame}

\begin{frame}[t,label=cuts]{Constraints from a failed run}
\small
\begin{tabularx}{\linewidth}{@{}lX@{}}
\toprule
Constraint & What it rules out\\
\midrule
Whole assignment & Only this complete vector of firm audits.\\[4pt]
Component assignment & Every plan keeping the failed component's firm decisions.\\[4pt]
Component prefix & Every plan preserving the relevant local exposure statuses.\\
\bottomrule
\end{tabularx}
\par\medskip
For the three-by-three failure, let $v_e$ indicate whether pair $e$ is visible.
A valid constraint is
\[
 (1-v_{ax})+(1-v_{by})+v_{cx}\geq1.
\]
At $(v_{ax},v_{by},v_{cx})=(1,1,0)$, the left side is zero: this pattern fails.
\sources{Changing a term is necessary, not sufficient. We choose the whole audit plan
before implementation, so a revised plan can omit an optional audit.}
\end{frame}

\begin{frame}[t,label=guarantees]{What the solver certifies}
\small
The master gives a lower bound $L$. A plan that passes the DA check gives
an upper bound $U$:
\[
 L\leq\resp\leq U.
\]
Every failure excludes a new audit pattern. With exact master solves,
the algorithm eventually certifies an optimum. A timeout leaves a verified
plan and a gap $U-L$.

\medskip
On the shared-firm construction, prefix cuts need at most two master solves.
Whole-assignment cuts need at least $2^N-1$ failed DA checks.
\sources{This round bound compares these cut methods. It is not a polynomial-time
guarantee. Proofs: Appendices 4--5. Numerical certification uses solver tolerances.}
\end{frame}
